% !TEX root = ../main.tex
\chapter{仕様書を Lean に移す方法}
\label{chap:spec-to-lean}
\BookIndex{けいしきしよう}{形式仕様}
\BookIndex{じゅつご}{述語}
\BookIndex{じぜんじょうけん}{事前条件}
\BookIndex{じごじょうけん}{事後条件}

\begin{chapterabstract}
自然言語で書かれた仕様書から、Lean で扱える小さな検証モデルを作る手順を学びます。
入力、出力、状態、操作、事前条件、事後条件、不変条件に分け、検査したい核を型、関数、述語、定理として切り出します。
題材には口座からの出金仕様を使い、成功時と失敗時の性質をモデル化します。
\end{chapterabstract}

\begin{learninggoals}
\item 自然言語仕様の曖昧な表現を、検査可能な部品へ分解できます。
\item 仕様の部品を、Lean の型、関数、述語、定理へ対応付けられます。
\item 事前条件、事後条件、不変条件の違いを説明できます。
\item Lean で検査する性質と、実装・運用で確認する項目を区別できます。
\item 口座出金仕様を、小さな Lean モデルとして表現できます。
\end{learninggoals}

\section{自然言語の仕様を Lean へ移す}

実務の仕様書は、人間どうしの合意を作るために書かれます。
したがって、読みやすさや業務文脈が重視されます。
一方、Lean が扱えるのは、型、関数、命題、証明として明示されたものだけです。
この差を埋めるには、仕様書をそのまま Lean に移さず、検査したい性質を小さく切り出す必要があります。

たとえば、仕様書に次のような文があったとします。

\begin{quote}
利用者は、残高の範囲内で口座から出金できる。残高を超える金額を出金しようとした場合は、エラーとして扱い、残高は変更しない。
\end{quote}

人間であれば、この文をおおよそ理解できます。
しかし Lean には、ここに出てくる「利用者」「口座」「残高」「範囲内」「出金できる」「エラー」「変更しない」が何なのか分かりません。
そこで、まず語を分けます。

\begin{itemize}
\item 口座は、少なくとも残高を持つデータです。
\item 出金額は自然数として扱います。
\item 出金可能とは、出金額が残高以下であることとします。
\item 出金に成功したら、残高は出金額だけ減ります。
\item 出金に失敗したら、状態は変わりません。
\end{itemize}

このように分けると、自然言語の仕様は少しずつ Lean の部品に近づきます。

\FigurePlaceholder{fig-ch24-spec-to-lean-pipeline.png}{0.92}{仕様の形式化パイプライン}{fig:ch24-spec-to-lean-pipeline}

図\ref{fig:ch24-spec-to-lean-pipeline}は、本章で使う分解の流れを示します。
左側には自然言語の仕様があります。
中央では、仕様を入力、状態、操作、条件に分けます。
右側では、それぞれを Lean の型、関数、述語、定理として表します。
下には、検証モデルと実装・運用上の確認項目の対応を置きます。
この対応により、「証明済み」という言葉が指す対象を明確にできます。

\section{自然言語仕様の危うさ}

自然言語仕様が危うい理由は、自然言語が悪いからではありません。
自然言語は、業務担当者、開発者、運用者、利用者の間で意味を共有するために必要です。
ただし、機械的な検査には向いていない表現があります。

典型的には、次のような表現が曖昧さを生みます。

\begin{itemize}
\item 「正常に処理する」：何をもって正常としますか。
\item 「適切なエラーを返す」：エラーの型、値、状態変化は何ですか。
\item 「一貫性を保つ」：どの値どうしの関係を保ちますか。
\item 「残高を更新する」：成功時だけですか。
失敗時にも何か記録しますか。
\item 「不正な入力を拒否する」：拒否した後の状態は変わりますか。
\end{itemize}

Lean に移すためには、これらの表現を、より小さな質問へ分解します。

本章では、仕様を次の六つの部品に分けます。
\begin{enumerate}
\item 入力：操作に渡される値です。
\item 出力：操作の戻り値です。
\item 状態：操作の前後で変化しうる値です。
\item 事前条件：操作を正しく呼ぶため、または成功するために必要な条件です。
\item 事後条件：操作後に満たすべき条件です。
\item 不変条件：操作の前後を通じて壊れてはいけない条件です。
\end{enumerate}

事前条件という語には、二つの使い方があります。
一つは、条件を満たさない入力を関数の定義域から除く使い方です。
もう一つは、全域関数の成功を保証する仮定として使う方法です。
本章の \texttt{withdraw} は任意の口座と金額を受け取る全域関数なので、\texttt{CanWithdraw} は後者の成功条件であり、呼び出し自体の適格性条件ではありません。

すべての仕様が、この六つへきれいに分かれるとは限りません。
むしろ、分けようとしたときに曖昧さが見つかります。
これは失敗ではありません。
Lean による形式化には、証明を書く前に仕様の穴を見つけるという価値があります。

\section{仕様を型・関数・述語・定理に分ける}

Lean へ移すときは、まず文を「名詞」「動詞」「条件」「主張」に分けるとよいでしょう。
名詞は型や構造体になりやすく、動詞は関数になりやすい要素です。
条件は述語になりやすく、主張は定理になりやすい要素です。

\begin{center}
\begin{tabular}{p{0.22\linewidth}p{0.34\linewidth}p{0.34\linewidth}}
\toprule
仕様書の要素 & 例 & Lean での表現 \\
\midrule
名詞 & 口座、残高、出金額 & 型、構造体、フィールド \\
動詞 & 出金する、更新する、拒否する & 関数 \\
条件 & 残高の範囲内、残高を超える & 述語、命題 \\
主張 & 成功時は残高が減る、失敗時は変わらない & 定理 \\
外部要素 & 外部決済、監査ログ、UI表示 & モデルとの対応を別途定義 \\
\bottomrule
\end{tabular}
\end{center}

この表は厳密な文法規則ではなく、設計レビューで使うための作業表です。
たとえば「残高を超える金額を出金しようとした場合はエラー」という文は、「条件」と「出力」と「状態保存」の三つを含んでいます。
これを一つの文章のまま扱うと曖昧ですが、三つに分ければ Lean で扱いやすくなります。

まず、口座と出金操作の最小モデルを作ります。

\begin{listing}[htbp]
\begin{minted}{lean4}
namespace Chapter24

-- 本章では、金額を Nat で単純化する。
-- 現実の通貨、小数、丸めは後で境界として明示する。
structure Account where
  balance : Nat
  deriving Repr, DecidableEq

-- 出金が成功するための条件。
def CanWithdraw (a : Account) (amount : Nat) : Prop :=
  amount <= a.balance

-- 成功すれば更新後の口座を返す。
-- 失敗すれば none を返す。
def withdraw (a : Account) (amount : Nat) : Option Account :=
  if amount <= a.balance then
    some { balance := a.balance - amount }
  else
    none

end Chapter24
\end{minted}
\caption{\texttt{Account} と \texttt{withdraw}}
\label{lst:ch24-account-model}
\end{listing}

この段階では、まだ「証明」はほとんどありません。
しかし、仕様はすでにかなり明確になっています。
\texttt{Account} は残高だけを持ちます。
\texttt{CanWithdraw} は出金可能条件を表します。
\texttt{withdraw} は成功か失敗を \texttt{Option Account} で返します。

ここには、設計判断も一つ入っています。
失敗時に元の口座を返すのではなく、\texttt{none} を返しています。
これは、失敗を戻り値で明示する設計です。
後で「失敗時は状態が変わらない」という性質を表したい場合、\texttt{withdraw} だけでは状態を返していないため、別の wrapper 関数を定義する必要があります。
このような設計判断が見えることも、形式化の価値です。

\section{事前条件・事後条件・不変条件}

仕様を Lean に移すときは、事前条件、事後条件、不変条件を混同しないことが重要です。
三つは似ていますが、読むタイミングが違います。

事前条件は、操作の前に仮定する条件です。
事後条件は、操作の後に期待する条件です。
不変条件は、操作の前後を通じて維持したい条件です。

口座出金の例では、次のように読めます。

\begin{itemize}
\item 事前条件：出金額が残高以下です。
\item 成功時の事後条件：更新後の残高は、元の残高から出金額を引いた値です。
\item 失敗時の事後条件：状態を変えないか、失敗として返します。
\item 不変条件：残高は負になりません。
\end{itemize}

ただし、本章のモデルでは残高を \texttt{Nat} で表しています。
\texttt{Nat} は自然数なので、負の残高は型の時点で表せません。
つまり「残高は負にならない」という条件の一部は、述語ではなく型によって保証されています。

型で排除できる不正状態と、定理として証明する不正状態を混同しないでください。
\texttt{Nat} を使えば負の値は表せません。
しかし、現実の金額型に丸めや符号付き整数を使うなら、同じ主張は自動では成り立ちません。
Lean のモデルが何を単純化しているかを、必ず書きます。

成功時の事後条件を、Lean の述語として書いてみます。

\begin{listing}[htbp]
\begin{minted}{lean4}
namespace Chapter24

-- 成功時の事後条件。
def PostSuccess
    (before : Account) (amount : Nat) (after : Account) : Prop :=
  after.balance = before.balance - amount

-- 事前条件があれば、withdraw は成功する。
theorem withdraw_success_when_allowed
    (a : Account) (amount : Nat)
    (h : CanWithdraw a amount) :
    withdraw a amount = some { balance := a.balance - amount } := by
  have hle : amount <= a.balance := by
    simpa [CanWithdraw] using h
  simp [withdraw, hle]

-- 成功時の戻り値は、PostSuccess を満たす。
theorem withdraw_success_post
    (a : Account) (amount : Nat)
    (h : CanWithdraw a amount) :
    Exists (fun b =>
      withdraw a amount = some b /\ PostSuccess a amount b) := by
  exact Exists.intro { balance := a.balance - amount }
    (And.intro (withdraw_success_when_allowed a amount h) rfl)

end Chapter24
\end{minted}
\caption{\texttt{withdraw\_success\_post}}
\label{lst:ch24-postcondition}
\end{listing}

\texttt{have hle : ... := ...} は、証明途中の補助事実に \texttt{hle} という名前を付けます。

\texttt{¬ P} は命題 \texttt{P} の否定を表します。

\texttt{withdraw\_success\_when\_allowed} は、「事前条件が満たされているなら成功する」という定理です。
\texttt{withdraw\_success\_post} は、「成功した結果は事後条件を満たす」という形に言い換えています。
実務の仕様書では、この二つが一文にまとめて書かれることがよくあります。
Lean では、分けて書くほうが読みやすくなります。

\section{証明すべき性質を選ぶ}

形式化で失敗しやすいのは、最初から仕様書全体を Lean に移そうとすることです。
実務では、壊れると困る性質を選ぶべきです。
特に、次のような性質は Lean に向いています。

\begin{itemize}
\item すべての入力について成り立つべき等式です。
\item 操作後にも保存される不変条件です。
\item 成功時と失敗時で分かれる事後条件です。
\item リファクタリング前後の意味保存です。
\item マイグレーションの round-trip 性質です。
\end{itemize}

外部システムや利用者との相互作用は、Lean の検証モデルと実装・運用上の検査を組み合わせて確認します。

\begin{itemize}
\item 外部 API が仕様どおりに応答することです。
\item DB のトランザクション分離レベルが実運用で満たされることです。
\item UI 上で利用者が誤操作しないことです。
\item 監査ログが規制要件を完全に満たすことです。
\item 実装コードと Lean モデルが常に同期していることです。
\end{itemize}

これらを Lean の対象に含める場合は、外部システムの仮定と実装との対応付けをモデルへ追加します。
検証対象は、小さな命題として明確に切り出せる核から選びます。

証明対象は、壊れたときの影響が大きく、小さな命題へ分けられる性質から選びます。

失敗時の仕様も、Lean で明示できます。
出金できない場合は \texttt{none} を返すという性質を、定理にします。

\begin{listing}[htbp]
\begin{minted}{lean4}
namespace Chapter24

-- 事前条件が成り立たなければ、withdraw は失敗する。
theorem withdraw_failure_when_not_allowed
    (a : Account) (amount : Nat)
    (h : Not (CanWithdraw a amount)) :
    withdraw a amount = none := by
  have hnot : ¬ amount <= a.balance := by
    simpa [CanWithdraw] using h
  simp [withdraw, hnot]

-- amount が 0 なら、どの口座でも出金後に同じ口座が返る。
theorem withdraw_zero (a : Account) :
    withdraw a 0 = some a := by
  cases a
  simp [withdraw]

end Chapter24
\end{minted}
\caption{\texttt{withdraw\_failure\_when\_not\_allowed}}
\label{lst:ch24-failure-spec}
\end{listing}

\texttt{withdraw\_failure\_when\_not\_allowed} は、失敗ケースを仕様から漏らさないための定理です。
\texttt{withdraw\_zero} は小さな境界条件です。
出金額が \texttt{0} の場合に結果がどうなるかを明示しています。

\section{検証モデルと実装の対応を明示する}

Lean の証明を実装へ適用するには、モデルで検査する性質と、実装・運用で検査する項目の対応を明示します。

「Lean で証明した」という記述には、対応する型、関数、仮定、定理を添えます。
外部システム、並行性、永続化、丸め、権限、監査要件は、それぞれに対応する実装・運用上の検査と組み合わせます。

口座モデルと実装・運用上の検査の対応は、次のように整理できます。

\begin{center}
\begin{tabular}{p{0.42\linewidth}p{0.48\linewidth}}
\toprule
Lean モデルで検査する項目 & 実装・運用で検査する項目 \\
\midrule
残高を \texttt{Nat} とする単純モデル & 通貨単位、小数、丸め、桁あふれ \\
出金成功条件と戻り値 & 認証、認可、口座所有者の確認 \\
成功時・失敗時の戻り値仕様 & DB 更新、ロック、トランザクション \\
小さな事後条件の証明 & 監査ログ、通知、外部決済連携 \\
Lean モデル上の性質 & 本番実装との完全一致 \\
\bottomrule
\end{tabular}
\end{center}

この表は、Lean の証明を実務で使うための対応表です。
どこまで保証したのかが明確であれば、レビューでも運用でも使いやすくなります。

\section{サンプル仕様：口座からの出金}

ここまでの内容を、仕様書の形から Lean の形へ通して見てみます。
まず、自然言語仕様を少し整えます。

\begin{quote}
口座は残高を持つ。出金要求は金額を持つ。出金額が残高以下であれば、出金は成功し、残高は出金額だけ減る。出金額が残高を超える場合、出金は失敗する。失敗した場合、状態は変更しない。検証モデルでは残高と出金額を自然数で表し、通貨単位、丸め、認可、DB 永続化、監査ログは実装・運用上の検査項目として分離する。
\end{quote}

この仕様は、すでに Lean へ移しやすい形になっています。
名詞、操作、条件、成功時、失敗時、実装上の検査項目が分かれているためです。

\FigurePlaceholder{fig-ch24-withdraw-scope.png}{0.92}{出金モデルの境界}{fig:ch24-withdraw-scope}

図\ref{fig:ch24-withdraw-scope}は、出金仕様のスコープを示します。
中央に Lean で扱う小さなモデルを置き、外側に実装・運用で検査する要素を置きます。
中央を小さく保つほど、定理は読みやすくなります。
外側を無視せず、境界として明示します。

失敗時に「状態は変更しない」と表したい場合、\texttt{withdraw} の戻り値だけでは表現しにくい面があります。
そこで、失敗時には元の口座を返す wrapper を定義します。

\begin{listing}[htbp]
\begin{minted}{lean4}
namespace Chapter24

-- 出金に失敗した場合は、元の口座をそのまま返す。
def withdrawOrKeep (a : Account) (amount : Nat) : Account :=
  match withdraw a amount with
  | some b => b
  | none => a

-- 失敗条件では、状態が変わらない。
theorem withdrawOrKeep_failure_keeps
    (a : Account) (amount : Nat)
    (h : Not (CanWithdraw a amount)) :
    withdrawOrKeep a amount = a := by
  have hnot : ¬ amount <= a.balance := by
    simpa [CanWithdraw] using h
  simp [withdrawOrKeep, withdraw, hnot]

-- 成功条件では、残高が amount だけ減る。
theorem withdrawOrKeep_success_balance
    (a : Account) (amount : Nat)
    (h : CanWithdraw a amount) :
    (withdrawOrKeep a amount).balance = a.balance - amount := by
  have hle : amount <= a.balance := by
    simpa [CanWithdraw] using h
  simp [withdrawOrKeep, withdraw, hle]

end Chapter24
\end{minted}
\caption{\texttt{withdrawOrKeep}}
\label{lst:ch24-withdraw-or-keep}
\end{listing}

ここで、\texttt{withdraw} と \texttt{withdrawOrKeep} は役割が違います。
\texttt{withdraw} は成功か失敗かを戻り値で表し、\texttt{withdrawOrKeep} は失敗を元の口座へ写す補助関数です。
後者は失敗情報を捨てるため、エラーを利用する API の代わりにはなりません。
また、この純粋関数の等式は、本番 DB で副作用が起きなかったことまでは証明しません。

\section{仕様レビューで使う形式}

実務では、Lean コードだけを示しても、仕様レビューとしては不十分な場合がよくあります。
自然言語、対応する Lean 定義、証明済みの定理、対象外の範囲を対応させると、読みやすくなります。

\begin{center}
\begin{tabular}{p{0.32\linewidth}p{0.28\linewidth}p{0.30\linewidth}}
\toprule
仕様項目 & Lean 側の名前 & レビューで確認すること \\
\midrule
口座は残高を持つ & \texttt{Account.balance} & 残高以外の状態を省いてよいか \\
出金可能条件 & \texttt{CanWithdraw} & 条件は \texttt{amount <= balance} で十分か \\
出金操作 & \texttt{withdraw} & 失敗を \texttt{none} で表してよいか \\
成功時の残高 & \texttt{withdraw\_success\_post} & 事後条件は業務仕様と合うか \\
失敗時の状態保存 & \texttt{withdrawOrKeep\_}\allowbreak\texttt{failure\_keeps} & 失敗時に副作用を持たないモデルでよいか \\
非対象範囲 & 境界表 & DB、認可、監査ログを別途検査するか \\
\bottomrule
\end{tabular}
\end{center}

この表は、Lean と仕様書をつなぐインデックスです。
Lean の定理名を仕様項目に対応させておくと、仕様変更時にどの定理を見直すべきかが分かりやすくなります。

定理名は仕様項目と対応させ、対象、条件、結論を追跡できるようにします。

\section{圏論の言葉で見直す}

本章の主題は Lean への移し方ですが、その背景には圏論の見方もあります。
仕様を型、関数、述語、定理へ分けると、ソフトウェア変更を「構造を保つ変換」として読みやすくなります。

\begin{itemize}
\item 型は、扱う対象の世界を決めます。
\item 関数は、対象から対象への変換を表します。
\item 述語は、対象や変換に課す制約を表します。
\item 定理は、変換が制約や意味を保存することを表します。
\end{itemize}

出金操作は、口座と金額から結果を作る変換です。
成功条件のもとでは、残高に関する事後条件を満たします。
失敗条件のもとでは、状態を保存します。
ここで見ているのは単なる計算結果ではなく、「どの性質が変換後にも残るか」です。
この保存性の見方は、マイグレーション、関手、自然変換、Pullback の各概念にも共通します。

\section{本章のまとめ}

自然言語仕様を Lean に移すときは、証明したい核を型、関数、述語、定理へ分け、入力、出力、状態、操作、事前条件、事後条件、不変条件に整理します。
証明対象は「重要かつ小さい」性質から選び、形式化しない性質とモデル境界も明示します。

本章では口座出金を例に、成功時と失敗時の仕様を表しました。
